\chapter{Productization: From Requirements to Controlled Ramp}
\label{ch:productization}
\label{ch:product-readiness}
\label{ch:validation}
\label{ch:reliability}
\label{ch:manufacturing}
\label{ch:networking}

This chapter shows how a technically sound optical design becomes a controlled
product. It is one lifecycle, not five separate interview subjects. Detail that
used to live in separate readiness, qualification, manufacturing, and
operations chapters is relocated: lifecycle tables and discipline definitions in
\Cref{app:reliability-reference}; manufacturing depth in
\Cref{app:manufacturing-reference}; link-operation and fabric survey in
\Cref{app:ai-fabric-context}. Failure analysis remains
\Cref{ch:failure-modes}.

\textit{Read first:} the spine below, the discipline table
(\Cref{tab:validation-jobs}), and the eight interview questions.

\textit{Deep dive:} relocated narratives behind the appendix pointers above.

\keyidea{Productization is the path from a product claim to a controlled ramp
with residual risk that you can name. Architecture margin, measurement
integrity, lifetime mechanisms, factory control, and bounded field exposure are
one argument. Memorizing EVT/DVT/PVT labels or long process catalogs is not the
point.}

\section{One lifecycle}
\label{sec:productization-spine}

Use this order when you reason about readiness. Calendar work may overlap. A
later gate cannot replace missing earlier evidence.

\begin{quote}
\noindent Define the product claim
$\rightarrow$ close the architecture
$\rightarrow$ characterize margins
$\rightarrow$ verify the design
$\rightarrow$ validate the system
$\rightarrow$ qualify lifetime risks
$\rightarrow$ validate the factory
$\rightarrow$ run a controlled pilot
$\rightarrow$ ramp and monitor.
\end{quote}

The eleven-step table follows in this chapter
(\Cref{tab:ladder,sec:product-readiness}). Phase-name translation and ownership
detail remain in the relocated narrative
(\Cref{sec:phase-labels,tab:readiness-ownership}).

\paragraph{Define the product claim.}
Write measurable requirements at named reference planes: optical, electrical,
thermal, management, reliability, and manufacturing. Name the release decision
each requirement supports. Without a claim, later data has no decision context.

\paragraph{Close the architecture.}
Ask whether optical, electrical, thermal, control, lifetime, and factory
budgets have a credible path. Catch thin margins before hardware makes changes
expensive. Derating and corner intent belong here
(\Cref{sec:derating-rules,sec:ladder-architecture}).

\paragraph{Characterize margins.}
Map nominal behavior, spread, sensitivities, and failure edges. Build a
behavioral model before you treat a number as a requirement result
(\Cref{sec:ladder-characterization}).

\paragraph{Verify the design.}
Compare evidence with frozen requirements at named planes and conditions.
Verification answers ``does it meet the contract?'' It does not by itself prove
intended-use fitness or lifetime.

\paragraph{Validate the system.}
Prove the product works in the intended host, peer, fiber, firmware, thermal,
and workload envelope. System validation is narrower than productization
(\Cref{tab:validation-jobs,sec:ladder-system-validation}).

\paragraph{Bring-up before characterization.}
Bring-up proves that hardware, firmware, management states, and a basic link
work reproducibly. Until that state exists, characterization numbers are not
interpretable. Name power rails, clocks, CMIS or equivalent state, optical
enable sequence, and the first BER or FEC observation that shows the link is
alive (\Cref{sec:bringup,sec:ladder-bringup}).

\section{Evidence disciplines in one table}

Keep the distinctions. Do not turn them into a dozen interview definitions.

\begin{table*}[ht]
\refstepcounter{table}\label{tab:validation-jobs}%
\centering
\setlength{\tabcolsep}{2.5pt}%
\footnotesize
\renewcommand{\arraystretch}{1.12}%
\begin{tabular}{@{}p{0.18\linewidth}p{0.26\linewidth}p{0.26\linewidth}p{0.24\linewidth}@{}}
\toprule
Discipline & Primary question & Main output & What it does not replace \\
\midrule
Characterization & How does the design behave and vary? & Behavioral model,
distributions, sensitivities, failure boundaries & Requirement verification or
release approval \\
\addlinespace[0.5em]
Verification & Does measured evidence meet the frozen requirement? & Traceable
requirement result with conditions and uncertainty & Intended-use evidence \\
\addlinespace[0.5em]
System validation & Is the complete product suitable for its intended system
use? & Supported operating and interoperability envelope & Lifetime
qualification or factory capability \\
\addlinespace[0.5em]
Reliability qualification & Will time and exposure cause unacceptable permanent
degradation? & Mechanism-specific life and environmental confidence & Every-unit
screening or manufacturing reproducibility \\
\addlinespace[0.5em]
Manufacturing validation & Can the production system repeatedly build, measure,
trace, and control the design? & Production capability, measurement confidence,
control plan, ramp evidence & System validation or life qualification \\
\addlinespace[0.5em]
Production acceptance testing & Should this unit or population proceed? & Pass,
fail, hold, rework, or disposition record & Complete mechanism coverage or
design qualification \\
\addlinespace[0.5em]
Pilot deployment & Does a bounded production-representative population behave as
predicted operationally? & Controlled field evidence and expansion decision &
Uncontrolled mass deployment \\
\addlinespace[0.5em]
Fleet monitoring & Does the deployed population match the release model? &
Trends, cohorts, alerts, incident evidence, next-revision learning & Root-cause
confirmation by itself \\
\bottomrule
\end{tabular}
\par\medskip
{\raggedright\footnotesize\textbf{Table~\thetable.} Evidence disciplines and
the decisions they can honestly support. Abbreviations:
\Cref{app:abbreviations}.\par}
\end{table*}

The same measurement can serve more than one discipline; the claim and decision
change, not the instrument.

Program names such as EVT, DVT, and PVT are company containers, not universal
proof categories. Interpret a phase gate by asking which build, which evidence,
which residual risks, and which decision
(\Cref{sec:phase-labels,tab:phase-readiness-map}).

\section{The optical product-readiness lifecycle}
\label{sec:product-readiness}
\label{sec:validation-workflow}

\refstepcounter{table}\label{tab:ladder}\label{tab:ladder-gates}%
\endgraf\medskip\noindent
\fcolorbox{keyidearule}{codebg}{%
\begin{minipage}{\dimexpr\textwidth-2\fboxsep-2\fboxrule\relax}
{\raggedright\footnotesize\textbf{Table~\thetable.} Optical product-readiness
lifecycle at a glance (Steps 1--11). Wall-chart:
\Cref{sec:tree-qualification}. Interview drill:
\Cref{sec:interview-validation-ladder}. NPI names: \Cref{tab:npi}.\par}
\medskip
\textbf{The product-readiness lifecycle removes a different uncertainty at each
step. The steps are logically ordered, although much of the work overlaps in
calendar time.}
{\raggedright\footnotesize
\begin{itemize}
\item \mbox{\textbf{Step 1: Define requirements.}} Uncertainty: what must the
      product do, and under what conditions? Decision: establish the evidence
      contract.
\item \mbox{\textbf{Step 2: Review architecture.}} Uncertainty: can the proposed
      design plausibly close the budgets and risks? Decision: proceed, modify,
      or reject the architecture.
\item \mbox{\textbf{Step 3: Bring up hardware.}} Uncertainty: is there a stable,
      known, reproducible operating state? Decision: begin interpretable
      characterization.
\item \mbox{\textbf{Step 4: Characterize behavior.}} Uncertainty: how does the
      design behave and vary? Decision: select corners and freeze evidence
      targets.
\item \mbox{\textbf{Step 5: Verify and validate.}} Uncertainty: does it meet
      requirements and work in the intended system? Decision: define the
      supported operating envelope. This step alone does not make the product
      ready for unrestricted production.
\item \mbox{\textbf{Step 6: Qualify reliability.}} Uncertainty: will time and
      exposure cause unacceptable permanent degradation? Decision: support the
      life and environmental claim.
\item \mbox{\textbf{Step 7: Validate manufacturing.}} Uncertainty: can
      production reproduce and control the design? Decision: authorize
      controlled production exposure.
\item \mbox{\textbf{Step 8: Run controlled pilot.}} Uncertainty: does a bounded
      deployed population behave as predicted? Decision: expand, hold,
      restrict, or roll back.
\item \mbox{\textbf{Step 9: Ramp mass production.}} Uncertainty: do factory,
      supplier, and field indicators remain controlled as exposure grows?
      Decision: increase production volume.
\item \mbox{\textbf{Step 10: Monitor fleet.}} Uncertainty: does deployed
      behavior match the release model? Decision: contain, sustain, or
      investigate.
\item \mbox{\textbf{Step 11: Feed learning forward.}} Uncertainty: which
      assumptions, controls, or architectures must change? Decision: improve
      the next revision.
\end{itemize}
\textit{Each step answers a question that the previous step cannot. A
later-stage success cannot compensate for missing evidence from an earlier
stage.}\par}
\end{minipage}%
}%
\endgraf\medskip

\section{Why program-phase names are not enough}
\label{sec:phase-labels}

EVT\sidenote{EVT: Engineering Validation Test or Engineering Verification
Test, depending on the organization.},
DVT\sidenote{DVT: Design Validation Test or Design Verification Test,
depending on the organization.},
and PVT\sidenote{PVT: Production Validation Test or Production Verification
Test, depending on the organization.}
are company containers for program timing, not universal technical
definitions. Sample size, exit criteria, and which evidence lands in which
phase vary by organization. Manufacturing-gate lookup for the same names lives
in \Cref{tab:npi}. Expanded EVT/DVT/PVT goal lists live in
\Cref{sec:phase-labels-detail}.

\keyidea{EVT, DVT, and PVT describe program timing and maturity.
Characterization, verification, system validation, reliability qualification,
and manufacturing validation describe what the evidence proves. When someone
says EVT, DVT, or PVT, ask four questions. Which hardware and software
configuration is included? Which engineering evidence is expected? Which risks
remain intentionally open? What decision or exit gate does the phase enable?}

After pilot, programs often speak of
MP\sidenote{MP: mass production. MP describes production exposure and
operating cadence, not evidence quality by itself.}
and sustaining. Those labels still need an exit decision.

\begin{table*}[ht]
\refstepcounter{table}\label{tab:phase-readiness-map}%
\centering
\setlength{\tabcolsep}{3pt}%
\footnotesize
\renewcommand{\arraystretch}{1.15}%
\begin{tabular}{@{}p{0.20\linewidth}p{0.28\linewidth}p{0.46\linewidth}@{}}
\toprule
Program label & Common emphasis & Typical product-readiness work \\
\midrule
EVT & Architecture and engineering learning & Architecture review, prototype
bring-up, risk retirement, early characterization, measurement development \\
\addlinespace[0.6em]
DVT & Design evidence and intended-use closure & Characterization completion,
requirement verification, margin, interoperability, system validation,
substantial qualification evidence \\
\addlinespace[0.6em]
PVT & Production-intent and ramp readiness & Production-reference freeze,
manufacturing validation, measurement-system analysis, ATP coverage, process
capability, controlled ramp evidence \\
\addlinespace[0.6em]
Pilot / limited deployment & Bounded operational evidence & Cohort deployment,
enhanced telemetry, success criteria, rollback, fleet comparison \\
\addlinespace[0.6em]
MP / sustaining & Controlled volume and field learning & Ramp, SPC, supplier
controls, fleet monitoring, failure analysis, next-revision feedback \\
\bottomrule
\end{tabular}
\par\medskip
{\raggedright\footnotesize\textbf{Table~\thetable.} Illustrative mapping from
program-phase labels to product-readiness work. EVT, DVT, and PVT are not
standards, and their activities overlap. Always ask what hardware population,
evidence, and exit decision the phase label represents. Reliability planning
may begin in EVT, representative qualification may run through DVT, and
corrective requalification may continue into PVT.\par}
\end{table*}



\paragraph{Who owns which evidence.}
\begin{table*}[ht]
\refstepcounter{table}\label{tab:readiness-ownership}%
\centering
\setlength{\tabcolsep}{3pt}%
\footnotesize
\begin{tabular}{@{}p{0.48\linewidth}p{0.46\linewidth}@{}}
\toprule
Topic & Primary owner \\
\midrule
IM/DD architecture and link budget & \Cref{ch:imdd} \\
Noise, sensitivity, RIN, BER, and metric interpretation & \Cref{ch:models} \\
Sources and modulation & \Cref{ch:lasers} \\
WDM and wavelength control & \Cref{ch:wdm} \\
Advanced packaging, 2.5D/3D, and CPO design judgment &
\Cref{ch:packaging} \\
AI/HPC rack architecture, fat-tree, rails, oversubscription &
\Cref{ch:hpc} \\
Product-readiness sequence and system-validation decisions &
\Cref{ch:product-readiness} \\
Reliability qualification & \Cref{ch:reliability} \\
Manufacturing validation and production control & \Cref{ch:manufacturing} \\
Optical-link operation, FEC, recovery, telemetry, and availability &
\Cref{ch:networking} \\
AI fabric context (topologies, module styles, CPO/XPO survey) &
\Cref{app:ai-fabric-context} \\
Failure analysis and recurrence control & \Cref{ch:failure-modes} \\
Instruments, procedures, and test-reference material &
\Cref{app:measurement-reference} \\
Decision navigation & \Cref{app:decision-trees} \\
\bottomrule
\end{tabular}
\par\medskip
{\raggedright\footnotesize\textbf{Table~\thetable.} Where detailed ownership
lives. The productization chapter keeps the readiness sequence and
decisions.\par}
\end{table*}



\section{Mechanism-driven qualification}
\label{sec:productization-qualification}

Reliability qualification is a bounded confidence argument, not a copied test
list.

\begin{quote}
\noindent Claim
$\rightarrow$ mechanism
$\rightarrow$ stress
$\rightarrow$ observable
$\rightarrow$ acceptance
$\rightarrow$ samples
$\rightarrow$ confidence
$\rightarrow$ decision
$\rightarrow$ handoff.
\end{quote}

The stress must accelerate the field mechanism without creating a different
one. Define observables and acceptance before the stress runs. Representative
samples matter; one convenient lot is not the released population. Zero
failures still leave a one-sided upper bound, not a zero rate
(\Cref{sec:qual-argument,sec:fit-dppm,sec:qual-sample-strategy}).

HTOL and similar stresses support life claims. Burn-in is an optional
production screen. Neither replaces the other
(\Cref{sec:qual-terms,app:reliability-reference}). Standards such as GR-468
supply methods and language; they are not the whole engineering argument
(\Cref{sec:qual-standards}).

\paragraph{Reversible heat versus permanent wear.}
A temperature sweep on a healthy unit asks how the product behaves while hot.
A qualification stress asks whether the exposure caused lasting change. Do not
treat a hot pass as proof that the same corner is safe after aging
(\Cref{tab:qual-activity-split}).

\paragraph{Failure timing.}
Early fails, random-in-time fails, and wear-out fails point to different
controls. Infant mortality may justify a screen. Wear-out needs a life model
and derating. Mixing them into one FIT number without stating the regime hides
the decision (\Cref{sec:bathtub,sec:wearout-modes}).

\paragraph{FIT and DPPM in one breath.}
One FIT is one failure per $10^9$ device-hours. DPPM counts defective units in
a named population and window. FIT uses time exposure; DPPM uses inspected
units. A FIT claim needs population, exposure, failure definition, model, and
confidence. Component FIT is not fabric availability
(\Cref{sec:fit-dppm,sec:fabric-availability}).

\paragraph{Acceleration validity.}
An Arrhenius or other acceleration model is only useful when the named
mechanism is active in the assumed regime and the stress did not invent a new
failure physics. If the stress creates a different mechanism, the result does
not support the field claim (\Cref{sec:accel-validity}).

\section{Factory evidence}
\label{sec:productization-manufacturing}

Manufacturing validation asks whether the production system can repeatedly
build, measure, trace, and control the released design.

Freeze the production reference first: design, BOM, process, firmware, test
software, fixtures, limits, and approved suppliers
(\Cref{sec:mfg-production-reference}). Build representative populations and keep
genealogy (\Cref{sec:mfg-builds-genealogy}).

Trust the measurement system before you interpret yield
(\Cref{sec:gauge-rr}). Separate first-pass yield from final yield: final yield
can hide rework and unstable process health
(\Cref{sec:yield-analysis,sec:yield-recovery-accounting}). Establish stability
before $C_p$/$C_{pk}$. Remember the three limit sets: control limits describe
process behavior, specification limits define product acceptance, and ATP
limits disposition units (\Cref{tab:mfg-three-limits,sec:process-capability}).

Validate that ATP catches known bad cases, not only good units
(\Cref{sec:hvm-test}). Every reaction plan needs a trigger, owner, containment,
evidence, restart rule, and follow-up (\Cref{tab:mfg-reaction-plan}). Ramp only
when evidence supports the next volume step
(\Cref{sec:mfg-controlled-ramp,app:manufacturing-reference}).

\paragraph{First-pass versus final yield.}
Suppose 1000 eligible units enter a route. Nine hundred pass on the first valid
attempt (FPY $=90\%$). After retest and approved rework, 990 are eventually
accepted (final yield $=99\%$). The nine-point gap is the recovery path. Track
it by failure mode, station, supplier, lot, and rework action. A release story
that quotes only 99\% hides the process health signal
(\Cref{sec:yield-recovery-accounting,tab:yield-split}).

\paragraph{Stability before capability.}
Plot the distribution, confirm the process is stable over time, then compute
$C_p$/$C_{pk}$ against the specification. An unstable process makes capability
numbers unreliable. Control limits are not specification limits; ATP limits are
a third set used to ship, hold, or reject
(\Cref{tab:mfg-three-limits,sec:process-capability}).

\paragraph{Representation before sample size.}
Sample the populations that actually vary: wafers or lots, assembly lines,
firmware revisions, and suppliers. Genealogy must survive rework. A large sample
of the wrong population still misleads
(\Cref{sec:mfg-builds-genealogy,sec:mfg-zero-defect-bound}).

\section{Pilot, ramp, and fleet feedback}
\label{sec:productization-pilot}

Link-up is not link health. A product can pass bring-up and still hurt a
workload through rising corrected-error demand, short bursts, or recovery
storms. The operational translation of physical margin lives in
\Cref{app:ai-fabric-context} and the relocated operations narrative
(\Cref{sec:ops-linkup-not-enough,sec:ops-what-system-sees}).

A pilot is a bounded field experiment: identifiable units, production-intent
configuration, enhanced telemetry, exit rules, and rollback
(\Cref{sec:ops-pilot-evidence,sec:ladder-fleet}). A small shipment without those
controls is not a pilot.

Expand only when metrics match the release model and no unexplained cohort
appears. Pause when risk is widening, unexplained, or hard to contain. Fleet
monitoring feeds escapes and margin discoveries back into requirements, ATP,
qualification, and the next revision
(\Cref{sec:ladder-feedback,sec:fabric-availability}).

\paragraph{What the system sees.}
Translate physical impairment into system language before you argue health.
Pre-FEC activity shows burden on the physical link. Corrected errors show FEC
working while margin may be shrinking. Uncorrectables and retrains are recovery
events with duration and repeat count. Average BER can hide short dangerous
bursts (\Cref{sec:ops-what-system-sees,sec:ops-bursts-transitions,sec:ops-recovery}).

\paragraph{Severity and cohorts.}
Operational severity follows workload consequence, not only optical symptoms.
Compare affected and unaffected groups that share lot, host, firmware, site, or
age. Correlation narrows ownership; it does not prove mechanism
(\Cref{sec:ops-severity,sec:ops-cohorts,ch:failure-modes}).

\paragraph{Availability is not FIT.}
Fabric availability also depends on event duration, redundancy, reroute, and
repair time. A rare per-link event can become frequent at fleet scale
(\Cref{sec:fabric-availability,sec:networking-availability}).

\section{What to protect when schedule is cut}

Protect, in order:
\begin{enumerate}
\item measurable requirements and named reference planes;
\item measurement integrity (otherwise later data lies);
\item thin architectural margins that cannot be recovered in the factory;
\item mechanism coverage for the dominant life risks;
\item factory controls that prevent shipping unknown populations;
\item a reversible pilot before uncontrolled exposure.
\end{enumerate}

Compress schedule by overlapping work and shrinking sample depth where residual
risk is already bounded. Do not skip the claim, the architecture close, or
measurement trust.

\section{Steps 1--11}
\label{sec:readiness-steps}

Use \Cref{tab:ladder} as the wall chart. The subsections below keep the step
labels and handoff rules. Derating examples, characterization maps, and Step~5
evidence-domain detail live in
\Cref{sec:readiness-steps-detail,sec:derating-rules}.

\subsection{Step 1: Define the requirements}
\label{sec:ladder-requirements}

Define performance, environment, reliability, manufacturing, and operational
requirement classes with owners and named planes
(\Cref{sec:systems-engineering-loop,tab:laser-prd}).

\textit{Evidence and handoff:} Freeze a signed requirements slice; refuse
hardware spend until success is defined.

\subsection{Step 2: Review the architecture}
\label{sec:ladder-architecture}

Close optical-power, noise, bandwidth, electrical, thermal, wavelength-control,
reliability, manufacturing, and serviceability budgets on stated assumptions
(\Cref{sec:systems-engineering-loop,tab:laser-engineering-checklist}). Apply
mechanism-based derating so temperature, aging, process spread, and measurement
uncertainty do not consume the last margin (\Cref{sec:derating-rules}).

\textit{Evidence and handoff:} Proceed to bring-up only when budgets close or
open items are named with redesign triggers.

\subsection{Step 3: Bring up the hardware}
\label{sec:ladder-bringup}

Establish a known, reproducible operating state. Separate integration fails
from product-performance questions (\Cref{sec:bringup}).

\textit{Evidence and handoff:} Continue when identity, supply rails,
management-ready state, basic optical power, lock, and a simple link are
reproducible.

\subsection{Step 4: Characterize the behavior}
\label{sec:ladder-characterization}

Map nominal behavior, distributions, sensitivities, interactions, failure
boundaries, and signatures. Averages can hide weak tails
(\Cref{sec:ladder-characterization-detail,sec:evidence-disciplines}).

\textit{Evidence and handoff:} Name thin ledgers and candidate corners; proceed
to Step~5, derate, or redesign before loaded-fleet claims.

\subsection{Step 5: Verify requirements and validate system use}
\label{sec:ladder-margin}
\label{sec:ladder-margin-interop}
\label{sec:ladder-interop}

Verification closes frozen requirements at named planes. System validation
closes intended use across supported hosts, peers, firmware, fiber plants, and
workloads (\Cref{sec:evidence-disciplines,sec:ladder-verify}). Track optical,
electrical, thermal, and control ledgers separately
(\Cref{sec:link-budget,tab:evidence-domains}). Golden-host margin is not the
ecosystem exit (\Cref{sec:prod-corners,sec:ladder-system-validation}).

\textit{Evidence and handoff:} Freeze requirement results, margin boundaries,
and supported combinations before qualification claims rely on them.

\subsection{Step 6: Qualify reliability}
\label{sec:ladder-stress}

Convert the operating envelope and credible mechanisms into a bounded life and
environmental confidence argument. This is not a harder system-validation
suite (\Cref{ch:reliability,sec:productization-qualification}).

\textit{Evidence and handoff:} Accept, derate, or hold life risk before
unrestricted ramp.

\subsection{Step 7: Validate manufacturing}
\label{sec:ladder-production}

Prove the factory can reproduce and control the supported design: production
reference, measurement system, yield, ATP, genealogy, and reaction plans
(\Cref{ch:manufacturing,sec:productization-manufacturing}).

\textit{Evidence and handoff:} Authorize controlled production exposure only
when multi-lot yield and escape detection support it.

\subsection{Step 8: Run a controlled pilot}
\label{sec:ladder-fleet}

Bounded production-representative deployment with cohort identity, success
criteria, enhanced telemetry, and rollback
(\Cref{sec:evidence-disciplines,ch:networking,sec:productization-pilot}).

\textit{Evidence and handoff:} Expand, hold, restrict, or roll back from exit
metrics versus the release model.

\subsection{Step 9: Ramp mass production}
\label{sec:ladder-mp}

Sustain volume under ATP, SPC, supplier, and change controls after pilot exit
(\Cref{ch:manufacturing,tab:npi}).

\textit{Evidence and handoff:} Increase volume only while factory, supplier, and
early-field indicators remain controlled.

\subsection{Step 10: Monitor the fleet}
\label{sec:ladder-fleet-monitor}

Compare deployed cohorts with the release model. Correlation prioritizes
hypotheses; it does not confirm a mechanism
(\Cref{sec:fec-symbol-histogram,sec:fleet-triage,ch:failure-modes}).

\textit{Evidence and handoff:} Contain, sustain, or investigate; return evidence
to the appropriate earlier readiness step when the release model fails.

\subsection{Step 11: Feed learning into the next revision}
\label{sec:ladder-feedback}

Convert fleet and manufacturing evidence into changed requirements, designs,
qualification, ATP, or controls (\Cref{sec:systems-engineering-loop}).

\textit{Evidence and handoff:} Write next-revision targets with owners, or
explicitly accept residual risk without change.

\section{Choosing evidence within a step}
\label{sec:ladder-evidence-class}

An experiment is not inherently characterization, verification, validation, or
qualification. Its category is determined by the claim being made
(\Cref{sec:evidence-disciplines}). Use the lowest-cost evidence that changes the
decision. State reference plane, condition, population, uncertainty, and
applicability. Reuse evidence across claims only when those conditions support
the new claim.

A first discriminating measurement should separate broad ownership classes; for
an optical-link symptom, named-plane power is often an inexpensive first split
between gross attenuation and signal-quality hypotheses
(\Cref{sec:validation-fork,sec:debug-fork}).

\heuristic{Name the reference plane before you name the instrument. A pretty
eye at the wrong plane is a wrong answer.}

\keyidea{Product readiness is a sequence of decisions, not a catalog of tests.
Each step should produce evidence for one explicit decision and identify the
uncertainty that remains.}

\section{Bring-up and production-representative corners}
\label{sec:bringup}

Bring-up proves a module and system can be powered, managed, and linked the way
production will run them. A quiet room-temperature BERT result is not system
validation.

\begin{enumerate}
\item Confirm identity, rails, temperature, and firmware.
\item Reach the required management-ready state.
\item Enable the intended lanes or light under host control.
\item Establish the optical path and verify basic power.
\item Establish electrical lock and lane alignment.
\item Run basic traffic and capture a reproducible baseline.
\end{enumerate}

A module forced into an emitting state and passing BER has not passed bring-up
if its required management sequence, safe state, diagnostics, alarms, and
recovery behavior are incorrect. That includes pluggable and external-laser
forms such as
ELSFP\sidenote{ELSFP: External Laser Small Form-Factor Pluggable. An
OIF-defined faceplate-pluggable CW laser module for CPO with a 24-pin card
edge, CMIS management, and blind-mate MT optics.}.
Register-level CMIS detail lives in
\Cref{sec:cmis,app:measurement-reference}.

\paragraph{Production-representative corners.}
\label{sec:prod-corners}

\begin{table*}[ht]
\refstepcounter{table}%
\centering
\setlength{\tabcolsep}{3pt}%
\footnotesize
\begin{tabular}{@{}p{0.28\linewidth}p{0.66\linewidth}@{}}
\toprule
Corner class & Representative challenges \\
\midrule
Thermal and loading & Case temperature, airflow, neighboring modules, full
traffic \\
\addlinespace[0.5em]
Electrical and host & Supported hosts, voltage corners, SerDes/equalization,
reset and restart \\
\addlinespace[0.5em]
Optical plant & Production fiber, connectors, reflections, loss, supported peers \\
\addlinespace[0.5em]
Control and service & CMIS transitions, alarms, hot-swap, recovery, firmware
combinations \\
\bottomrule
\end{tabular}
\par\medskip
{\raggedright\footnotesize\textbf{Table~\thetable.}\label{tab:prod-corners}
Production-representative corners for system validation. Mechanism detail:
\Cref{ch:lasers,ch:wdm,ch:networking}.\par}
\end{table*}


\section{Worked example: 800G DR4-class module}
\label{sec:productization-worked}

\textit{Claim.} Ship an 800G DR4-class module for a named host family, reach,
temperature range, pre-FEC BER objective, and FIT envelope.

\textit{Architecture close.} Confirm modulator and laser headroom, host
equalization assumptions, thermal path, CMIS state machine, and factory probe
points. Flag any budget that closes only at nominal.

\textit{Characterize.} Map Tx OMA, TDECQ, RIN under stated ORL, Rx sensitivity
at the named BER plane, and corner sensitivities across voltage and
temperature.

\textit{Verify and validate.} Trace requirements to named planes. Run
supported-host interoperability and loaded-chassis corners
(\Cref{sec:prod-corners}).

\textit{Qualify.} Pick mechanisms (laser wear-out, solder fatigue, moisture,
electrostatic damage). Choose stresses that accelerate those mechanisms. State
sample basis and residual risk (\Cref{sec:qual-worked-laser}).

\textit{Factory.} Freeze BOM and firmware. Clear gauge R\&R on the ATP
stations that gate shipment. Track first-pass yield by failure mode. Set
reaction plans for BER and extinction-ratio fallout.

\textit{Pilot.} Deploy a bounded cohort with enhanced FEC and thermal
telemetry, hold criteria, and rollback ownership. Expand only if distributions
match the release model.

\textit{Monitor.} Watch corrected-error and retrain cohorts by lot, host, and
site. Feed escapes into ATP and the next revision
(\Cref{sec:ops-worked-retrains,ch:failure-modes}).

\section{Interview takeaway}
\label{sec:productization-design-review}

\keyidea{I show that the architecture can become a controlled product: claim,
margin evidence, life argument, factory control, and reversible field exposure.
I do not recite an NPI org chart. When evidence conflicts, I escalate the claim
that is least supported, not the loudest dashboard.}

\section{Interview Q\&A}
\label{sec:productization-interview-qa}

Practice aloud. Prefer first-person reasoning. Score with
\Cref{sec:chapter-interview-rubric}. Extended productization drills live in the
secondary bank (\Cref{sec:cheat-flashcards,app:reliability-reference,app:manufacturing-reference}).

\paragraph{Question 1. How do you know a new optical design is ready to move
forward?}
\textit{Tests:} architecture margin; evidence; residual risk; decision-making.

\textit{Spoken answer.}
``I ask what claim we are trying to support next, what evidence closes the
thin margins, and what residual risk remains. If the architecture still depends
on an unproven budget, or if the measurement plane is undefined, I do not call
it ready. Ready means the next gate has a named decision, not that every test
in the lab passed once.''

\textit{Pressure follow-up.} ``We have good bench BER. Why wait?''\\
\textit{Answer pivot.} ``Bench BER without corners, host coverage, life
mechanisms, or factory measurement trust is not a release argument.''

\textit{Trap:} ``Ready means the prototype meets the datasheet at room
temperature.''

\paragraph{Question 2. How do characterization and system validation differ?}
\textit{Tests:} behavioral model versus intended-use evidence.

\textit{Spoken answer.}
``Characterization maps how the design behaves and where it fails.
System validation asks whether the complete product works in the intended
system and workload. I can characterize deeply on a reference host and still
fail system validation on production hosts''
(\Cref{tab:validation-jobs}).

\textit{Pressure follow-up.} ``Where does verification fit?''\\
\textit{Answer pivot.} ``Verification checks a frozen requirement at a named
plane. It is neither a full behavioral map nor a full intended-use proof.''

\textit{Trap:} ``Characterization, verification, and validation are just
increasingly hard test levels.''

\paragraph{Question 3. How do you build a mechanism-driven reliability plan?}
\textit{Tests:} claim, mechanism, acceleration, observable, acceptance.

\textit{Spoken answer.}
``I start from the life or environment claim, name the mechanism that could
break it, choose a stress that accelerates that mechanism, define the
observable and acceptance rule before stressing, then size samples for the
confidence I need. Zero fails still leave an upper bound''
(\Cref{sec:qual-argument,sec:fit-dppm}).

\textit{Pressure follow-up.} ``Can we skip HTOL if burn-in looks clean?''\\
\textit{Answer pivot.} ``Burn-in is a screen. It does not replace the
qualification claim unless the argument explicitly says so.''

\textit{Trap:} ``We passed GR-468, so reliability is done.''

\paragraph{Question 4. How do you know the factory can reproduce the design?}
\textit{Tests:} measurement trust; yield; capability; controls.

\textit{Spoken answer.}
``I freeze the production reference, clear the measurement system, read
first-pass yield and failure Pareto, confirm stability before capability, and
check that ATP and reaction plans catch real defects. Final yield alone is not
enough'' (\Cref{sec:gauge-rr,sec:yield-analysis,tab:mfg-three-limits}).

\textit{Pressure follow-up.} ``Final yield is 99\%. Ship?''\\
\textit{Answer pivot.} ``Not until I know first-pass yield and what recovery is
costing and hiding.''

\textit{Trap:} ``A passing FAIR means the factory is ready at volume.''

\paragraph{Question 5. What makes a pilot useful?}
\textit{Tests:} bounded exposure; observability; rollback.

\textit{Spoken answer.}
``A useful pilot has a hypothesis, production-representative hardware,
enhanced telemetry, exit rules, and rollback. Expansion requires metrics that
match the release model and no unexplained cohort''
(\Cref{sec:ops-pilot-evidence}).

\textit{Pressure follow-up.} ``We already shipped a small lot. Is that a
pilot?''\\
\textit{Answer pivot.} ``Not if exposure, telemetry, and rollback were never
defined.''

\textit{Trap:} ``Any early customer shipment counts as a pilot.''

\paragraph{Question 6. Schedule is cut. What evidence do you protect?}
\textit{Tests:} risk-based prioritization.

\textit{Spoken answer.}
``I protect requirements and planes, measurement integrity, thin architectural
margins, dominant life mechanisms, factory controls that prevent unknown
populations, and a reversible pilot. I compress calendar overlap before I erase
those.''

\textit{Pressure follow-up.} ``Management wants to skip the pilot.''\\
\textit{Answer pivot.} ``Then the residual risk acceptance must be explicit,
owned, and reversible some other way. Skipping silently is not a plan.''

\textit{Trap:} ``Cut whatever is last on the schedule.''

\paragraph{Question 7. Supplier data says pass, but system behavior says
otherwise. What do you do?}
\textit{Tests:} component versus system evidence.

\textit{Spoken answer.}
``I treat both as evidence under different claims. I re-check planes,
conditions, populations, and whether the supplier test answers the system
failure mode. I do not discard system evidence because a component ATP passed,
and I do not blame the supplier until ownership is localized''
(\Cref{sec:supplier-exec,ch:failure-modes}).

\textit{Pressure follow-up.} ``Can we ship while we investigate?''\\
\textit{Answer pivot.} ``Only with bounded exposure, enhanced monitoring, and
an owned containment trigger.''

\textit{Trap:} ``Supplier pass clears the system fail.''

\paragraph{Question 8. Give a 60-second productization plan.}
\textit{Tests:} end-to-end command of the spine.

\textit{Spoken answer.}
``Define the claim and planes, close the architecture, characterize margins,
verify requirements, validate intended use, qualify the dominant life
mechanisms, prove the factory can measure and control the design, run a
bounded pilot with rollback, then ramp while monitoring cohorts and feeding
escapes back. At each gate I name the decision, the evidence, and the residual
risk.''

\textit{Pressure follow-up.} ``Where do you start if hardware already exists?''\\
\textit{Answer pivot.} ``I still write the claim and freeze the measurement
planes first. Otherwise the existing data cannot support a release decision.''

\textit{Trap:} ``Start wherever the lab is already busy.''
